\chapter{Experiment \& Results}
\label{chap:3}
% =========================================================
% Chapter 4 -- Sections 4.1 and 4.2
% Ready-to-paste LaTeX
%
% Suggested appendix placeholders to add later:
%   app:metrics:definitions      -> formal metric definitions (PDR, delay, success rate, etc.)
%   app:stats:aggregation        -> mean, standard deviation, coefficient of variation
%   app:stats:correlation        -> Pearson r, Spearman rho, p-values
%   app:stats:windowing          -> time-window aggregation and filtering rules
%   app:stats:normalisation      -> rate normalisation by run duration / per-run aggregation
% =========================================================

\section{Experimental Design and Evaluation Methodology}
\label{sec:results:methodology}
To keep the main text focused on interpretation, formal metric definitions and
post-processing rules are reported in Appendix~\ref{app:metrics:definitions},
Appendix~\ref{app:stats:aggregation}, and Appendix~\ref{app:stats:windowing}.
\subsection{Experimental Objective and Research-Question Mapping}
\label{sec:results:rq-mapping}

The purpose of the experiments is to determine how different UAV charging
scheduling policies affect system performance when charging coordination is carried
over the same FANET used for telemetry. Since all other scenario conditions are
held constant, the comparison isolates the effect of the scheduling policy on both
resource allocation and network behaviour.

The evaluation is organised around the three research questions introduced in
Chapter~1:

\begin{itemize}
    \item \textbf{RQ1 -- FANET robustness and routing quality.}
    The first objective is to assess whether different charging outcomes are associated
    with different levels of telemetry reliability, measured primarily through packet
    delivery ratio (PDR), end-to-end delay, and their temporal evolution.

    \item \textbf{RQ2 -- Charging policy, mission continuity, and survivability.}
    The second objective is to compare the ability of each policy to preserve UAV
    availability through successful charging, reduced waiting burden, and lower
    battery depletion rates.

    \item \textbf{RQ3 -- Shared communication substrate and cross-layer coupling.}
    The third objective is to examine whether degradation in the FANET also
    weakens the charging-control loop itself, thereby creating a feedback mechanism
    in which poor connectivity leads to failed charging, further depletion, and
    additional network degradation.
\end{itemize}

\subsection{Protocols Under Comparison}
\label{sec:results:protocols}

Seven charging policies are evaluated in Test Round~2. They include both
non-preemptive and preemptive strategies:

\begin{itemize}
    \item \texttt{ugv\_fcfs}: non-preemptive first-come, first-served;
    \item \texttt{ugv\_edf}: non-preemptive earliest deadline first;
    \item \texttt{ugv\_role\_priority}: non-preemptive role-priority scheduling;
    \item \texttt{ugv\_dynamic}: non-preemptive dynamic-score scheduling;
    \item \texttt{ugv\_p\_edf}: preemptive earliest deadline first;
    \item \texttt{ugv\_p\_role\_priority}: preemptive role-priority scheduling;
    \item \texttt{ugv\_p\_dynamic\_score}: preemptive dynamic-score scheduling.
\end{itemize}

The non-preemptive policies decide only which waiting UAV should be served next
when a dock becomes available. By contrast, the preemptive policies are allowed to
interrupt a lower-priority charging allocation when a more urgent request arrives.

\subsection{Replication Strategy, Dataset Scope, and Artifact Availability}
\label{sec:results:dataset-scope}

Each of the seven protocols was executed in three independent replicate runs,
yielding a total of 21 runs for Test Round~2. All runs span the same simulated
mission duration of 180 minutes and were produced under the same scenario
configuration apart from the policy switch. This balanced design enables both
within-policy stability checks and cross-policy comparisons under matched
experimental conditions.

In the overall experimental workflow, \textbf{Test Round~1} was used primarily as a
preliminary validation stage to assess simulator stability, inspect logging behaviour,
and identify implementation issues before conducting the final comparative study.
The results presented in this chapter are instead based on \textbf{Test Round~2},
which is the main evaluation dataset used for protocol comparison.

A further important difference is that the weather node in Test Round~2 was
configured in \textbf{Markov mode}, so that weather regimes evolve over time through
macrostate transitions rather than remaining fixed throughout the entire run. This
choice makes the evaluation more representative of time-varying environmental
stress and allows the charging and network subsystems to be tested under changing
conditions.

The analysis therefore adopts the protocol as the primary comparison factor and the
replicate run as the unit of repeated observation. Aggregate protocol summaries are
reported from the set of three runs belonging to the same policy, while temporal
plots preserve run-level trajectories in order to expose instability, outliers, and
failure cascades that would be hidden by averaging alone.
The chapter concludes with a controlled follow-up validation under fixed sunny
weather, designed to test whether the main survivability-related pattern observed in
Test Round~2 remains visible when weather variance is reduced.

For reproducibility, the experiment outputs and the plotting/analysis scripts used to
generate the figures discussed in this chapter are maintained in the project GitHub
repository.\footnote{\url{https://github.com/kk-Syuer/UAV_UGV_monitoring}} In
particular, the Test Round~2 dataset is organised under
\texttt{experiment\_data\_collection/test\_round2}, while the corresponding figures are
stored under \texttt{analysis/test\_round2/figures/}.

% \subsection{Logged Data Sources Used in the Analysis}
% \label{sec:results:data-sources}

% The results in this chapter are derived from the append-only logs introduced in
% Section~3.9. In particular, the analysis relies on a subset of the files that directly
% capture network quality, charging decisions, charging execution, and UAV survival.

% The most important files are the following:

% \begin{itemize}
%     \item \texttt{qos\_metrics.csv}, which provides cumulative reliability and delay
%     statistics for telemetry and control flows;

%     \item \texttt{network\_timeseries.csv}, which records time-windowed network
%     indicators such as window-level PDR and mean delay, enabling temporal analysis
%     of degradation episodes;

%     \item \texttt{charge\_events.csv}, which records the end-to-end outcome of each
%     charging request, including success, timeout, routing-drop cases, and decision-side
%     timing/context fields when available;

%     \item \texttt{charge\_session\_events.csv}, which records dock-level charging
%     activity such as session start, completion, and charged energy;

%     \item \texttt{charge\_queue\_timeseries.csv}, which captures queue length and
%     dock occupancy over time;

%     \item \texttt{death\_events.csv}, which records battery depletion events and
%     therefore provides the basis for survivability analysis;

%     \item \texttt{weather\_timeseries.csv}, which provides the environmental context
%     used to interpret weather-driven degradation and routing-drop episodes.
% \end{itemize}

\subsection{Evaluation Metrics}
\label{sec:results:metrics}

The evaluation uses four families of metrics, each corresponding to one aspect of the
system behaviour under study.

\paragraph{Network-quality metrics.}
The primary network indicator is the packet delivery ratio (PDR), used as the main
measure of telemetry reliability. End-to-end delay is also reported, but interpreted
with caution because later sections show that raw delay values may be biased by
delivery conditioning when the network collapses and only a restricted subset of
packets continues to arrive. Window-level PDR and delay are further used to study
temporal degradation and recovery.

\paragraph{Charging-performance metrics.}
Charging effectiveness is evaluated using charging success rate, timeout rate,
ROUTING\_DROP rate, decision latency, effective wait, per-session recovered energy,
and dock utilisation. These metrics distinguish between failures due to queue
pressure and failures due to the inability of the request to reach the scheduler.

\paragraph{Fleet-survivability metrics.}
Mission continuity is assessed through the number and temporal distribution of
battery depletion events. The analysis also distinguishes between depletions affecting
cluster heads and member UAVs, since the loss of a cluster head carries a larger
topological consequence for the FANET backbone.

\paragraph{Cross-layer metrics.}
To support the main argument of the chapter, network and charging indicators are
not analysed in isolation. Instead, the results later examine statistical associations
between PDR and selected charging-related variables, including charge success rate,
ROUTING\_DROP rate, dock utilisation, queue length, and depletion burden. The
purpose is not to claim strict causality from correlation alone, but to test whether the
observed protocol differences are consistent with the proposed mechanism linking
charging coordination to network quality. Formal definitions of these metrics are
provided in Appendix~\ref{app:metrics:definitions}.

\subsection{Dataset Reliability Note}
\label{sec:results:data-reliability}

Before comparing protocol outcomes, the Test Round~2 dataset was checked for
completeness and structural consistency. All seven protocols were available with
three replicate runs each, the required analysis files were present for all runs, and
all runs reached the target mission duration. Additional preprocessing notes,
including sentinel-value handling, duplicated windows, and record-filtering rules, are
reported in Appendix~\ref{app:data-validation}.

\section{Aggregate Protocol-Level Results}
\label{sec:results:aggregate}

This section compares the seven charging policies at an aggregate level using the
three replicate runs available for each protocol. The discussion focuses on three
dimensions: overall network performance, charging effectiveness, and fleet
survivability. Detailed metric definitions and aggregation rules are reported in
Appendix~\ref{app:metrics:definitions} and Appendix~\ref{app:stats:aggregation}.

\subsection{Overall Network Performance}
\label{sec:results:network-kpis}

At the aggregate level, \texttt{ugv\_p\_edf} achieves the best overall network
performance, with the highest mean packet delivery ratio among all tested policies.
\texttt{ugv\_fcfs} and \texttt{ugv\_dynamic} follow, while
\texttt{ugv\_role\_priority} and \texttt{ugv\_p\_role\_priority} form the weakest
tail of the ranking, as shown in Fig.~\ref{fig:mean-pdr-merged}.
A second relevant observation is the difference in stability across replicates.
Although the mean PDR range across protocols is not extremely wide, some
policies show much higher between-run variance than others. In particular,
\texttt{ugv\_edf} and \texttt{ugv\_p\_role\_priority} display visibly larger
dispersion, indicating that their behaviour is less robust under the same nominal
conditions. This instability becomes more visible in the time-series overlays of
Fig.~\ref{fig:pdr-events-variance-all} and becomes important later when analysing
failure cascades and outlier runs.
Delay is treated as a secondary indicator throughout; its interpretation is revisited in
Section~\ref{sec:results:delay-bias}.

Overall, the aggregate network comparison already suggests that preemptive,
urgency-aware scheduling is beneficial at system level. The advantage of
\texttt{ugv\_p\_edf} is not limited to queue handling alone; rather, it appears
associated with better preservation of the operational fleet and therefore with more
reliable multi-hop communication.

\begin{figure}[h]
  \centering
  \includegraphics[width=.8\textwidth]{experiment_plots/01_mean_pdr_merged.pdf}
  \caption{Mean packet delivery ratio by protocol in Test Round 2 under stochastic
(Markov-mode) weather from \texttt{qos\_metrics.csv}. Data source: Test Round 2 (Markov-regime weather).}
  \label{fig:mean-pdr-merged}
\end{figure}


\subsection{Charging Scheduler Performance}
\label{sec:results:charging-kpis}

The charging results reinforce the network-level ranking. Among the seven
protocols, \texttt{ugv\_p\_edf} achieves the highest charge success rate, indicating
that preemptive deadline-based scheduling is the most effective at converting
requests into successful dock starts under contention. \texttt{ugv\_dynamic} is also
strong, combining a high success rate with comparatively low decision latency and
effective wait (Fig.~\ref{fig:charge-success-rate}).

At the opposite end, \texttt{ugv\_role\_priority} and \texttt{ugv\_edf} show the
lowest success rates. Their weaker performance is reflected not only in the number
of successful charging outcomes, but also in the lower energy recovered per session
and the higher prevalence of unresolved requests. This indicates that poor charging
performance is not merely a matter of slower queue progression, but a more general
difficulty in sustaining the service loop under adverse conditions.

An important result is that routing-related request loss remains a major failure mode
across all protocols. In other words, a non-negligible share of charging failures
occurs before the scheduler can even make a decision, as also shown by the
outcome breakdown in Fig.~\ref{fig:charge-outcome-breakdown}. This supports the thesis
claim that the charging subsystem is tightly coupled to network health: when the
FANET is degraded, some UAVs are unable to reliably reach the UGV with their
requests, and scheduling quality alone cannot compensate for that loss.

For this reason, the charging metrics should be interpreted jointly. A protocol with a
moderate decision time but a high success rate is more valuable than one with fast
decisions on the small subset of requests that actually arrive. Later sections return
to this distinction and show that request reachability is more consequential than raw
decision latency.  


\begin{figure}[h]
  \centering
  \includegraphics[width=.8\textwidth]{experiment_plots/02_charge_success_rate_merged.pdf}
  \caption{Charging success rate by protocol. Data source: Test Round 2 (Markov-regime weather).}
  \label{fig:charge-success-rate}
\end{figure}

\begin{figure}[h]
  \centering
  \includegraphics[width=.8\textwidth]{experiment_plots/03_charge_outcome_breakdown_merged.pdf}
  \caption{Breakdown of charging outcomes by protocol. Data source: Test Round 2 (Markov-regime weather).}
  \label{fig:charge-outcome-breakdown}
\end{figure}

\subsection{Fleet Survivability}
\label{sec:results:survivability}

Fleet survivability further clarifies the protocol differences. The lowest depletion
burden is observed under \texttt{ugv\_p\_edf}, followed by \texttt{ugv\_dynamic},
while \texttt{ugv\_role\_priority} and especially \texttt{ugv\_edf} show markedly
worse outcomes. This confirms that the most effective charging policies are also the
ones that best preserve the availability of the aerial fleet over the full mission
horizon, as shown in Fig.~\ref{fig:dead-uav-cumulative}.

The depletion statistics are particularly important because UAV loss has a direct
network consequence. When member UAVs deplete, sensing and forwarding capacity
are reduced; when cluster heads deplete, the impact is even more severe because
the communication backbone itself becomes unstable. For this reason, depletion
count is not only a survivability metric but also a mechanism-level indicator that
helps explain subsequent differences in PDR.

The aggregate results also reveal a difference between structural weakness and
outlier-driven instability. \texttt{ugv\_role\_priority} performs poorly in a relatively
consistent way across runs, which suggests an intrinsic limitation of the policy.
By contrast, \texttt{ugv\_edf} is strongly affected by one extreme run, indicating that
its average result is partly shaped by catastrophic instability rather than uniformly
poor behaviour in every replicate.

Taken together, the aggregate evidence points to a coherent pattern:
protocols that achieve more successful charging outcomes tend to show fewer
depletions, and protocols with fewer depletions tend to preserve better network
reliability.

\begin{figure}[h]
  \centering
  \includegraphics[width=.8\textwidth]{experiment_plots/03_dead_uav_cumulative_merged.pdf}
  \caption{Cumulative UAV battery depletion events by protocol. Data source: Test Round 2 (Markov-regime weather).}
  \label{fig:dead-uav-cumulative}
\end{figure}

\subsection{Aggregate Trade-Off View}
\label{sec:results:tradeoff-view}

When the metrics are read together, \texttt{ugv\_p\_edf} emerges as the strongest
overall policy in Test Round~2. It combines the highest mean PDR, the highest
charging success rate, and the lowest depletion burden, making it the most balanced
solution among the evaluated protocols. \texttt{ugv\_dynamic} is also competitive,
especially in terms of scheduler responsiveness, but its overall network and
survivability performance remains slightly weaker.

At the opposite end, \texttt{ugv\_role\_priority} is the weakest policy overall. Its
lower charging success, higher timeout burden, and heavier depletion load are
consistent with its weaker network outcomes. \texttt{ugv\_p\_role\_priority} is also
unfavourable, although for partly different reasons, since its poor network ranking
coexists with a comparatively healthier surviving fleet and therefore requires more
careful interpretation in the delay and cross-layer analysis.

These aggregate results establish the broad ranking of the policies, but they do not
yet explain \emph{how} degradation develops over time. For that reason, the next
section moves from protocol-level averages to temporal trajectories of PDR, queue
pressure, dock utilisation, and depletion events.

\begin{figure}[h]
  \centering
  \includegraphics[width=.8\textwidth]{experiment_plots/04_policy_radar_merged.pdf}
  \caption{Aggregate trade-off across major protocol KPIs. Data source: Test Round 2 (Markov-regime weather).}
  \label{fig:policy-radar}
\end{figure}

\section{Temporal Dynamics of Charging and Network Quality}
\label{sec:results:temporal}

\subsection{PDR Evolution Over Time}
\label{sec:results:pdr-time}

The time-series plots show that the protocols differ not only in average PDR, but
also in the way reliability evolves over the 180-minute mission. In general,
\texttt{ugv\_p\_edf} exhibits the most stable trajectory, with smaller and less
persistent drops in delivery performance. By contrast, \texttt{ugv\_edf} and
\texttt{ugv\_p\_role\_priority} show much larger between-run variability, including
episodes of abrupt degradation that are not visible from the mean alone.

This temporal view is important because some protocols achieve similar aggregate
PDR values while behaving very differently over time. A protocol with modest but
stable reliability may be preferable to one with a similar mean but repeated collapse
intervals. For this reason, the time-series plots in Fig.~\ref{fig:network-pdr-time-merged}
and Fig.~\ref{fig:pdr-events-variance-all} are used throughout the chapter to
separate structural weakness from outlier-driven behaviour. The dip-threshold rule
used in the event-aligned plots is defined in Appendix~\ref{app:stats:dip-threshold}.


\begin{figure}[h]
  \centering
  \includegraphics[width=.88\textwidth]{experiment_plots/01_network_pdr_over_time_merged.pdf}
  \caption{Merged PDR over time across protocols. Data source: Test Round 2 (Markov-regime weather).}
  \label{fig:network-pdr-time-merged}
\end{figure}

\begin{figure}[h]
  \centering
  \includegraphics[width=.88\textwidth]{experiment_plots/07_pdr_events_merged_variance_all.pdf}
  \caption{PDR evolution with replicate variance and event overlays. It reports the temporal mean of \texttt{window\_pdr} samples from \texttt{network\_timeseries.csv} (windowed 10-second sliding values, sentinel values excluded). Data source: Test Round 2 (Markov-regime weather).}
  \label{fig:pdr-events-variance-all}
\end{figure}

\subsection{Queue Pressure, Dock Utilisation, and Charging Throughput}
\label{sec:results:queue-dock}

The charging time-series help clarify whether poor outcomes are mainly caused by
persistent queue overload or by failures earlier in the control path. Across
protocols, dock utilisation remains within a relatively narrow range compared with
the much larger variation observed in timeout rates and routing-drop rates. This
suggests that unsuccessful charging is not explained solely by full docks or long
local queues.

A more informative pattern is the co-evolution of queue activity and charging
throughput. In healthier runs, requests continue to arrive, docks remain active, and
cumulative charged energy grows steadily. In weaker runs, charging throughput
flattens over time and dock utilisation may fall rather than rise, indicating that the
system is no longer sustaining a healthy flow of successful requests. In this sense,
low utilisation is not necessarily a sign of low demand; it can also be a symptom of
network degradation and fleet attrition.

This distinction is useful for interpreting timeout behaviour. If timeout burden grows
without a corresponding sustained increase in dock pressure, the cause is more
likely to lie in policy decisions or request reachability than in simple charger
saturation. These patterns are visible in Fig.~\ref{fig:queue-length-merged},
Fig.~\ref{fig:dock-util-merged}, and Fig.~\ref{fig:cumulative-energy-merged}.


\begin{figure}[h]
  \centering
  \includegraphics[width=.88\textwidth]{experiment_plots/03_charge_queue_length_merged.pdf}
  \caption{Queue length over time across protocols. Data source: Test Round 2 (Markov-regime weather).}
  \label{fig:queue-length-merged}
\end{figure}

\begin{figure}[h]
  \centering
  \includegraphics[width=.88\textwidth]{experiment_plots/03_dock_utilization_merged.pdf}
  \caption{Dock utilisation over time across protocols. Data source: Test Round 2 (Markov-regime weather).}
  \label{fig:dock-util-merged}
\end{figure}

\begin{figure}[h]
  \centering
  \includegraphics[width=.88\textwidth]{experiment_plots/03_cumulative_energy_charged_merged.pdf}
  \caption{Cumulative charged energy over time across protocols. Data source: Test Round 2 (Markov-regime weather).}
  \label{fig:cumulative-energy-merged}
\end{figure}

\subsection{Weather Context and Time-Varying Degradation}
\label{sec:results:weather-time}

Because Test Round~2 uses the weather node in Markov mode without fixed weather seed, network conditions are
subject to time-varying environmental stress rather than a fixed regime for the
entire experiment. The weather-conditioned plots therefore provide important
context for interpreting temporal performance changes. In particular, lower PDR and
higher delay are expected during harsher regimes, and these periods can coincide
with higher charging disruption.

The value of these plots is to show
that protocols are being evaluated under changing external conditions rather than a
fully static environment. This makes the comparison more realistic and also helps
explain why some degradation episodes appear clustered in time instead of being
uniformly distributed across the mission horizon.

For this reason, temporal drops in PDR or charging throughput should be read
together with the weather timeline in Fig.~\ref{fig:weather-timeseries} and with the
weather-conditioned comparisons in Fig.~\ref{fig:pdr-weather-merged} and
Fig.~\ref{fig:delay-weather-merged}, especially before drawing conclusions about
whether a failure is policy-driven, environment-driven, or the result of both factors
acting together.

\begin{figure}[h]
  \centering
  \includegraphics[width=.82\textwidth]{experiment_plots/05_pdr_vs_weather_regime_merged.pdf}
  \caption{Mean PDR by protocol and weather regime. Data source: Test Round 2 (Markov-regime weather).}
  \label{fig:pdr-weather-merged}
\end{figure}

\begin{figure}[h]
  \centering
  \includegraphics[width=.82\textwidth]{experiment_plots/05_delay_vs_weather_regime_merged.pdf}
  \caption{Mean delay by protocol and weather regime. Data source: Test Round 2 (Markov-regime weather).}
  \label{fig:delay-weather-merged}
\end{figure}

\begin{figure}[h]
  \centering
  \includegraphics[width=.88\textwidth]{experiment_plots/05_weather_timeseries.pdf}
  \caption{Weather regime evolution over time. Data source: Test Round 2 (Markov-regime weather).}
  \label{fig:weather-timeseries}
\end{figure}

\section{Cross-Layer Analysis: From Scheduling to Network Quality}
\label{sec:results:cross-layer}

The previous sections showed that the protocol ranking is broadly consistent across
network, charging, and survivability metrics. This section makes that relationship
explicit by analysing how charging outcomes are associated with FANET quality. The
goal is not to claim causality from correlation alone, but to test whether the observed
protocol differences are consistent with the mechanism proposed in this thesis:
better charging continuity preserves more UAVs, and a healthier fleet preserves more
stable multi-hop communication.

\subsection{Statistical Association Between Charging Outcomes and PDR}
\label{sec:results:cross-layer-corr}

Across the 21 runs, charge success rate is positively associated with mean PDR,
while charge-request \texttt{ROUTING\_DROP} rate is negatively associated with it.
In the round-2 analysis, charge success shows both positive Pearson and Spearman
correlations with PDR, whereas \texttt{ROUTING\_DROP} shows negative
correlations of similar magnitude. This indicates that protocols which convert more
requests into successful charging sessions also tend to maintain better telemetry
delivery, as shown in Fig.~\ref{fig:cl-charge-success-vs-pdr} and
Fig.~\ref{fig:corr-heatmap-pdr}.

A second important result is that battery depletion burden is strongly negatively
associated with PDR. This makes depletion a useful mediating variable between the
charging layer and the communication layer: poor charging outcomes are reflected in
more frequent UAV loss, and heavier UAV loss is reflected in lower network
reliability. In this sense, depletion is not only an outcome metric, but also a bridge
that helps explain why charging policy differences reappear at network level. 

\begin{figure}[h]
  \centering
  \includegraphics[width=.82\textwidth]{experiment_plots/CL_B_charging_vs_pdr_scatter.pdf}
  \caption{Association between charging success and mean PDR. Data source: Test Round 2 (Markov-regime weather).}
  \label{fig:cl-charge-success-vs-pdr}
\end{figure}

\begin{figure}[h]
  \centering
  \includegraphics[width=.82\textwidth]{experiment_plots/06_corr_heatmap_pdr.pdf}
  \caption{Correlation structure of PDR and selected charging indicators. Data source: Test Round 2 (Markov-regime weather).}
  \label{fig:corr-heatmap-pdr}
\end{figure}

\subsection{Mechanism Interpretation}
\label{sec:results:cross-layer-mechanism}

The aggregate and temporal results support a coherent mechanism chain. First, the
scheduling policy affects whether a charging request succeeds, times out, or fails
before reaching the scheduler. Second, unsuccessful charging reduces the amount of
energy effectively returned to the fleet. Third, reduced energy recovery increases the
probability of UAV depletion. Finally, heavier depletion weakens the relay structure
of the FANET, especially when cluster heads are affected, and this lowers packet
delivery performance. This chain is summarised in Fig.~\ref{fig:cl-mechanism-chain}
and its temporal signature is visible in Fig.~\ref{fig:cl-pdr-vs-depletions}.

This interpretation also explains why the strongest protocol in this study is
\texttt{ugv\_p\_edf}. Its advantage is not limited to queue efficiency at the charger.
Rather, it appears to preserve fleet continuity more effectively, and that continuity in
turn supports higher PDR. Conversely, the weakest protocols do not fail only because
they make poor local scheduling decisions; they also allow a larger degradation loop
to form between charging failure, UAV loss, and network deterioration.

\begin{figure}[h]
  \centering
  \includegraphics[width=.85\textwidth]{experiment_plots/CL_D_mechanism_chain.pdf}
  \caption{Cross-layer mechanism linking charging outcomes, depletion, and PDR. Data source: Test Round 2 (Markov-regime weather).}
  \label{fig:cl-mechanism-chain}
\end{figure}

\begin{figure}[h]
  \centering
  \includegraphics[width=.85\textwidth]{experiment_plots/CL_F_pdr_vs_depletions_timeseries.pdf}
      \caption{Temporal co-evolution of mean of \texttt{window\_pdr} samples from \texttt{network\_timeseries.csv} and depletion burden. Data source: Test Round 2 (Markov-regime weather).}
  \label{fig:cl-pdr-vs-depletions}
\end{figure}

\subsection{Why Request Reachability Matters More Than Decision Speed}
\label{sec:results:cross-layer-reachability}

One of the clearest findings of the cross-layer analysis is that decision latency and
effective waiting time are much weaker predictors of PDR than request reachability.
The round-2 report finds no strong or significant association between mean PDR and
raw decision latency, whereas \texttt{ROUTING\_DROP} rate shows a clear negative
association. In practical terms, this means that the dominant bottleneck is often not
how quickly the scheduler responds, but whether the request reaches the scheduler
at all.

This distinction is important for the interpretation of protocol quality. A policy may
appear computationally responsive once a request is available locally at the UGV, yet
still perform poorly at system level if many requests are lost in transit during adverse
network conditions. In the present architecture, charging coordination depends on
the same FANET that carries telemetry, so network degradation can directly disrupt
access to the charging service itself. 

The result therefore strengthens one of the main thesis claims: in a shared-substrate
design, charging performance cannot be evaluated only in terms of queue discipline
or local scheduling speed. It must also be evaluated in terms of end-to-end control
reachability under network stress, as shown in Fig.~\ref{fig:cl-correlation-bar} and
Fig.~\ref{fig:cl-routing-drop-time}.


\begin{figure}[h]
  \centering
  \includegraphics[width=.82\textwidth]{experiment_plots/CL_E_correlation_bar_chart.pdf}
  \caption{Correlation strength of selected charging indicators with PDR. Data source: Test Round 2 (Markov-regime weather).}
  \label{fig:cl-correlation-bar}
\end{figure}

\begin{figure}[h]
  \centering
  \includegraphics[width=.82\textwidth]{experiment_plots/CL_G_routing_drop_timeseries.pdf}
  \caption{Routing-drop burden over time. Data source: Test Round 2 (Markov-regime weather).}
  \label{fig:cl-routing-drop-time}
\end{figure}

\subsection{The CH Priority Paradox: Why Role-Based Protocols Underperform}
\label{sec:results:ch-priority-paradox}

The role-based protocols, \texttt{ugv\_role\_priority} and
\texttt{ugv\_p\_role\_priority}, perform worse than expected despite explicitly giving
cluster heads (CHs) the highest scheduling priority. In the present simulator, this
underperformance is not mainly caused by dock-level contention. Instead, it arises
from a more structural effect: the system contains only two CHs, and a CH disables
packet forwarding while charging. As a result, when both CHs are admitted to
charging in close temporal proximity, the relay backbone becomes temporarily
unavailable for the three member UAVs.

\paragraph{Hypothesis.}
The two CHs have identical battery capacity and broadly similar relay-intensive duty
cycles, so they tend to approach the charging threshold at nearly the same time.
This produces repeated paired CH charge requests. Under role-based scheduling,
both CH requests are promoted ahead of member requests and are likely to be served
in close succession, or even in parallel given the multi-dock UGV configuration.
Because a charging CH does not forward packets, simultaneous CH charging removes
the only CH-level relays available to the members. The members then lose their path
to the UGV scheduler, causing higher \texttt{ROUTING\_DROP} and lower charging
success. As member survivability deteriorates, the network weakens further and the
failure loop amplifies.

\paragraph{Empirical evidence.}
The left panel of Figure~\ref{fig:cl-ch-sync-cascade} shows that the two CHs often
generate charge requests within a short temporal gap across protocols, confirming
that CH demand is naturally synchronized. This is consistent with the fleet design:
the CHs have the same battery capacity and similar relay workload, so they deplete
on correlated time scales.

Figure~\ref{fig:cl-role-audit} shows that the role-based protocols are associated with
worse member-side outcomes than the stronger urgency-based alternatives,
including lower member charge success and heavier member routing-drop burden.
This pattern is consistent with the loss of relay availability when CH charging events
become temporally clustered. The protocol-level KPI overview in
Figure~\ref{fig:cl-protocol-kpi-overview} supports the same interpretation: despite
prioritising CHs, the role-based policies still exhibit weaker PDR, poorer overall
charge success, and heavier depletion burden.

The important point is that the paradox does not require CHs to wait behind each
other at the dock. Since the UGV offers multiple docking spots, both CHs can be
served quickly. However, this local scheduling success becomes a system-level
problem because it removes both relay-capable CHs from the forwarding layer at the
same time.

\paragraph{Interpretation and design implication.}
The main failure mode of role-based scheduling is therefore simultaneous CH
withdrawal from the relay backbone rather than simple queue inefficiency. In a fleet
with only two CHs, prioritising both of them too aggressively can momentarily
improve CH access to charging while severely damaging communication access for
the members.

A more suitable fix would be to retain role awareness but avoid concurrent CH
charging whenever both CHs are the only active high-level relays. One possible
design is a relay-preserving constraint: if one CH is already charging, the second CH
may be delayed briefly or conditionally admitted only when member connectivity can
still be preserved. This would reduce backbone collapse without abandoning the role
hierarchy entirely.

Further implementation details for the role-stratified metrics, CH synchronization
analysis, and relay-loss interpretation should be documented in the appendix
(Appendix~\ref{app:role-stratified-analysis},
Appendix~\ref{app:ch-sync-analysis}).
\begin{figure}[h]
    \centering
    \includegraphics[width=0.92\textwidth]{experiment_plots/CL_M_role_scheduling_audit.pdf}
    \caption{Role-stratified scheduling audit. In the role-based protocols, CHs do
    not receive genuinely effective priority, while members experience lower charging
    success and higher routing-drop burden. Data source: Test Round 2 (Markov-regime weather).}
    \label{fig:cl-role-audit}
\end{figure}

\begin{figure}[h]
    \centering
    \includegraphics[width=0.92\textwidth]{experiment_plots/CL_N_ch_sync_cascade.pdf}
    \caption{CH synchronization and member-starvation cascade. Left: CH charge
    requests occur in close temporal proximity across protocols. Right: runs with
    longer CH wait tend to exhibit weaker member charge success. Data source: Test Round 2 (Markov-regime weather).}
    \label{fig:cl-ch-sync-cascade}
\end{figure}

\begin{figure}[h]
    \centering
    \includegraphics[width=0.92\textwidth]{experiment_plots/CL_A_protocol_kpi_overview.pdf}
    \caption{Cross-protocol KPI overview. The role-based protocols remain weak in
    overall PDR, depletion burden, and charge success despite explicitly prioritising
    cluster heads. Data source: Test Round 2 (Markov-regime weather).}
    \label{fig:cl-protocol-kpi-overview}
\end{figure}
\section{Delay Interpretation and Measurement Bias}
\label{sec:results:delay-bias}

End-to-end delay must be interpreted alongside PDR, not independently. Delay is computed only over delivered packets; under heavy UAV attrition, many multi-hop transmissions are lost, leaving a surviving packet set that traverses a simpler, shorter topology. This produces a conditioning bias: protocols with worse survivability can appear to have lower delay precisely because their network has collapsed. The weak relationship between PDR and mean delay in Fig.~\ref{fig:pdr-vs-delay} confirms that delay carries limited protocol-discriminating information in this dataset. For this reason, the comparative evaluation is based on PDR, charging success, and depletion burden, while PDR-weighted delay is discussed only as a robustness check.
\begin{figure}[h]
  \centering
  \includegraphics[width=.82\textwidth]{experiment_plots/06_mean_e2e_delay_merged.pdf}
  \caption{Mean end-to-end delay by protocol. Data source: Test Round 2 (Markov-regime weather).}
  \label{fig:mean-delay-merged}
\end{figure}

\begin{figure}[h]
  \centering
  \includegraphics[width=.82\textwidth]{experiment_plots/06_pdr_vs_e2e_scatter.pdf}
  \caption{Joint view of PDR and mean end-to-end delay. Data source: Test Round 2 (Markov-regime weather).}
  \label{fig:pdr-vs-delay}
\end{figure}
The exact definition of the bias-corrected delay metric used in the robustness check
is provided in Appendix~\ref{app:stats:delay-robustness}; the base delay definition
remains in Appendix~\ref{app:metrics:definitions}.

The conditioning effect is illustrated by Fig.~\ref{fig:cl-delay-conditioning}, while
Fig.~\ref{fig:cl-delay-robustness} shows the corresponding robustness check.

\begin{figure}[h]
  \centering
  \includegraphics[width=.82\textwidth]{experiment_plots/CL_C_e2e_delay_conditioning_analysis.pdf}
  \caption{Conditioning-bias analysis for end-to-end delay. Data source: Test Round 2 (Markov-regime weather).}
  \label{fig:cl-delay-conditioning}
\end{figure}

\begin{figure}[h]
  \centering
  \includegraphics[width=.82\textwidth]{experiment_plots/CL_K_delay_robustness_panel.pdf}
  \caption{Robustness checks for delay interpretation. Data source: Test Round 2 (Markov-regime weather).}
  \label{fig:cl-delay-robustness}
\end{figure}
\section{Controlled Validation Under Fixed Weather}
% =============================================================================
% Section 4.6 — Controlled Validation Under Fixed Weather
% Drop-in LaTeX file for Chapter 4 of the thesis.
% =============================================================================

\subsection{Controlled Validation Under Fixed Weather}
\label{sec:results:controlled-validation}

The preceding sections established that charging scheduling policy,
fleet survivability, and network quality form a coherent cross-layer
mechanism under the stochastic weather conditions of Test Round~2.
A natural question is whether the observed trends are primarily driven
by weather variance itself or whether they reflect a more fundamental
property of the scheduling--survivability--network chain.
To answer this, Test Round~3 repeats the protocol comparison under
\emph{fixed sunny weather}, eliminating the Markov regime transitions
that introduced time-varying routing stress in Round~2.
If the cross-layer relationships persist despite the removal of weather
variance, the mechanism can be considered robust to the environmental
factor that dominated Round~2 degradation episodes.

% -----------------------------------------------------------------------------
\subsubsection{Rationale and Experimental Setup}
\label{sec:results:controlled-validation:setup}

\paragraph{Design changes from Round~2 to Round~3.}
The only deliberate change between the two rounds is the weather
configuration.  In Round~2, the weather node operated in Markov mode,
cycling through sunny, cloudy, rainy, and stormy regimes according to
stochastic transition probabilities.  In Round~3, the weather node is
locked to the \texttt{sunny} regime for the entire mission duration.
All \texttt{weather\_timeseries.csv} files in the Round~3 dataset
confirm a single regime label (\texttt{sunny}) throughout every run.
By fixing weather, WEATHER\_DROP-induced routing failures are removed,
which should reduce both ROUTING\_DROP rates and PDR variance across
all protocols.

\paragraph{What remains constant.}
The protocol set (seven policies), fleet composition (two cluster heads,
three member UAVs, one UGV with charging dock), simulated mission
duration ($\sim$180~min per run), number of replicates (three per
protocol), logging schema, and metric definitions are identical to
Round~2.  The analysis pipeline is also the same, ensuring that any
observed difference is attributable to weather control rather than
methodological change.

\paragraph{Dataset scope and filtering.}
The Round~3 dataset contains 25~run folders in total.  Four folders
whose names contain \texttt{stormy} (case-insensitive) are excluded
because they belong to a separate sensitivity check and violate the
fixed-sunny design.  After filtering, 21~sunny runs remain:
7~protocols $\times$ 3~replicates each, matching the Round~2 structure
exactly.

\paragraph{Threats to validity.}
Even with fixed weather, residual stochasticity remains from random
seeds governing packet generation timing, UAV mobility jitter, and
charging-request arrival patterns.  The validation therefore targets
\emph{trend consistency} across rounds---whether the same directional
relationships hold---rather than exact numerical reproduction of
Round~2 values.

% -----------------------------------------------------------------------------
\subsubsection{Results and Comparison with Test Round~2}
\label{sec:results:controlled-validation:results}

This section compares the Round~2 and Round~3 outcomes category by
category, using a structured \emph{trend match / mismatch} framework.
For each category, the Round~2 trend is stated first, followed by
the Round~3 observation, the match verdict, and the supporting figure
references.  A consolidated evidence map is provided at the end of
this subsection (Table~\ref{tab:evidence-map}).

% ---- 1. Network Reliability ------------------------------------------------
\paragraph{1) Network reliability: overall PDR and PDR over time.}

\emph{Round~2 trend.}
\texttt{ugv\_p\_edf} achieves the highest mean PDR (0.630), followed by
\texttt{ugv\_fcfs} (0.616) and \texttt{ugv\_dynamic} (0.595).  The
role-based protocols form the weakest tail, with
\texttt{ugv\_p\_role\_priority} at 0.554.  PDR time-series show high
between-run variance, particularly for \texttt{ugv\_edf} and
\texttt{ugv\_p\_role\_priority}, with episodic drops linked to weather
regime transitions.
(Figures: \texttt{01\_mean\_pdr\_merged.png},
\texttt{01\_network\_pdr\_over\_time\_merged.png},
\texttt{CL\_A\_protocol\_kpi\_overview.png} in
\texttt{test\_round2/figures/}.)

\emph{Round~3 observation.}
All protocols show an absolute PDR increase of 0.07--0.20, consistent
with the removal of weather-induced packet loss.
\texttt{ugv\_p\_role\_priority} rises to rank~1 (0.754),
\texttt{ugv\_p\_edf} drops to rank~2 (0.732), and
\texttt{ugv\_dynamic} falls to rank~7 (0.667).
The Kendall rank correlation between Round~2 and Round~3 PDR rankings
is $\tau = -0.24$, indicating no preservation of the ordinal ranking.
(Figures: \texttt{01\_mean\_pdr\_merged.png},
\texttt{01\_network\_pdr\_over\_time\_merged.png},
\texttt{CL\_A\_protocol\_kpi\_overview.png} in
\texttt{test\_round3/figures/}.)

\emph{Trend match:} \textbf{NO} for protocol ranking.
However, the positive association between charging success and PDR
is preserved at a qualitative level: the protocol with the highest
Round~3 charge success rate (\texttt{ugv\_p\_role\_priority}, 0.598)
is also the Round~3 PDR leader.

% ---- 2. Charging Outcomes ---------------------------------------------------
\paragraph{2) Charging outcomes: success / timeout / routing-drop
composition.}

\emph{Round~2 trend.}
\texttt{ugv\_p\_edf} leads in charge success rate (48.6\%), while
\texttt{ugv\_role\_priority} trails at 38.4\%.  ROUTING\_DROP rates
range from 0.218 (\texttt{ugv\_p\_edf}) to 0.346 (\texttt{ugv\_edf}),
and are a significant negative predictor of PDR ($r = -0.53$).
(Figures: \texttt{02\_charge\_success\_rate\_merged.png},
\texttt{03\_charge\_outcome\_breakdown\_merged.png},
\texttt{CL\_G\_routing\_drop\_timeseries.png} in
\texttt{test\_round2/figures/}.)

\emph{Round~3 observation.}
All protocols improve in success rate under sunny weather, with
\texttt{ugv\_p\_role\_priority} reaching the top (0.598) and
\texttt{ugv\_edf} recovering strongly to rank~2 (0.583).
ROUTING\_DROP rates collapse dramatically to 0.010--0.020 across all
protocols, confirming that WEATHER\_DROP was the dominant cause of
request loss in Round~2.
Timeout rates, however, increase substantially for all protocols
(e.g.\ \texttt{ugv\_dynamic} from 0.230 to 0.434), suggesting that
with longer-surviving fleets and fewer routing drops, queue contention
replaces routing failure as the primary bottleneck.
Kendall $\tau$ for charge success rate ranking: 0.33 (NO match).
(Figures: \texttt{02\_charge\_success\_rate\_merged.png},
\texttt{03\_charge\_outcome\_breakdown\_merged.png},
\texttt{CL\_G\_routing\_drop\_timeseries.png} in
\texttt{test\_round3/figures/}.)

\emph{Trend match:} \textbf{NO} for ranking, but the \emph{structural
finding} that routing drops are weather-coupled is strongly confirmed
by their near-elimination in Round~3.

% ---- 3. Survivability -------------------------------------------------------
\paragraph{3) Survivability: battery depletions.}

\emph{Round~2 trend.}
\texttt{ugv\_p\_edf} has the fewest depletions (20.0 per run),
\texttt{ugv\_edf} the most (45.0).  Higher depletion burden is
negatively associated with PDR ($r = -0.57$).
(Figures: \texttt{03\_dead\_uav\_cumulative\_merged.png},
\texttt{CL\_F\_pdr\_vs\_depletions\_timeseries.png} in
\texttt{test\_round2/figures/}.)

\emph{Round~3 observation.}
Depletions decrease for all protocols, with the most dramatic
reductions in the protocols that suffered worst in Round~2:
\texttt{ugv\_edf} drops from 45.0 to 16.7 (rank~1 in Round~3),
and \texttt{ugv\_p\_role\_priority} from 26.7 to 17.3 (rank~2).
The Kendall $\tau$ is $-0.05$ (NO ranking match).
(Figures: \texttt{03\_dead\_uav\_cumulative\_merged.png},
\texttt{CL\_F\_pdr\_vs\_depletions\_timeseries.png} in
\texttt{test\_round3/figures/}.)

\emph{Trend match:} \textbf{NO} for ranking.  The convergence of
depletion counts across protocols is consistent with the removal of
weather-driven cascade failures that disproportionately affected weaker
protocols in Round~2.  With routing drops near zero, even protocols
that previously lost many requests can now successfully charge,
compressing the survivability gap.

% ---- 4. Cross-Layer Associations --------------------------------------------
\paragraph{4) Cross-layer associations.}

\emph{4a) Success rate vs.\ PDR.}
In Round~2, charge success rate shows a positive association with PDR
(Pearson $r = 0.58$, $p = 0.006$).  In Round~3, the protocol with
the highest success rate (\texttt{ugv\_p\_role\_priority}) is also the
PDR leader.  However, because protocol rankings have shifted,
statistical confirmation requires recomputation of the per-run
correlation on the Round~3 dataset.
(Figures: \texttt{CL\_B\_charging\_vs\_pdr\_scatter.png} in both
rounds; \texttt{CL\_E\_correlation\_bar\_chart.png}.)

\begin{quote}\textbf{TODO:} Recompute Pearson and Spearman correlation
between charge success rate and mean PDR for Round~3 runs ($n = 21$).
Verify whether the positive association remains significant.  If the
correlation is computed in
\texttt{CL\_E\_correlation\_bar\_chart.png} for Round~3, read the
values from that figure and report them here.\end{quote}

\emph{4b) Routing-drop rate vs.\ PDR.}
In Round~2, ROUTING\_DROP rate is a strong negative predictor of PDR
($r = -0.53$).  In Round~3, routing drops collapse to near-zero for
all protocols, so the predictor effectively loses variance.  This
result does not contradict the Round~2 finding; rather, it confirms
that once routing drops are removed (by fixing weather), their
explanatory power over PDR disappears because the mechanism is no
longer active.
(Figures: \texttt{CL\_G\_routing\_drop\_timeseries.png},
\texttt{CL\_E\_correlation\_bar\_chart.png} in both rounds.)

\emph{4c) Depletions vs.\ PDR.}
In Round~2, depletions are negatively associated with PDR
($r = -0.57$).  In Round~3, the compressed depletion range makes the
statistical test less powerful, but the qualitative pattern persists:
protocols with fewer depletions (\texttt{ugv\_edf}, 16.7;
\texttt{ugv\_p\_role\_priority}, 17.3) achieve above-average PDR.
(Figures: \texttt{CL\_F\_pdr\_vs\_depletions\_timeseries.png},
\texttt{CL\_D\_mechanism\_chain.png} in both rounds.)

\emph{Trend match for cross-layer associations:} \textbf{PARTIAL.}
The directional relationships are qualitatively preserved, but the
compressed variance in Round~3 (due to weather control) weakens the
statistical power of the correlations.  The routing-drop pathway is
confirmed by its near-elimination.

% ---- 5. Delay Discussion ----------------------------------------------------
\paragraph{5) Delay discussion (cautious).}

End-to-end delay is computed only on \emph{delivered} packets.  As
discussed in Section~4.5, this introduces a conditioning bias: when
PDR drops, the surviving packet set traverses a sparser, shorter
topology, which can paradoxically \emph{reduce} measured delay.  This
caveat remains relevant in Round~3, although the bias should be
attenuated because PDR is uniformly higher and more stable under
fixed sunny weather.

In Round~2, mean E2E delay ranges from 52.8\,ms
(\texttt{ugv\_p\_dynamic\_score}) to 71.4\,ms
(\texttt{ugv\_p\_role\_priority}).  In Round~3, the range narrows to
48.1--61.6\,ms, with \texttt{ugv\_p\_edf} achieving the lowest delay
(48.1\,ms).  The Kendall $\tau$ for delay ranking is $-0.33$ (NO
match), and the ranking reversal---\texttt{ugv\_p\_role\_priority}
improves from worst to second-best---is consistent with the
conditioning bias: in Round~2, \texttt{ugv\_p\_role\_priority} had
the lowest PDR, so its delivered packets were biased toward short
paths; in Round~3, with much higher PDR, its delay rises less because
the full multi-hop topology is preserved.
(Figures: \texttt{06\_mean\_e2e\_delay\_merged.png},
\texttt{06\_pdr\_vs\_e2e\_scatter.png},
\texttt{CL\_C\_e2e\_delay\_conditioning\_analysis.png},
\texttt{CL\_K\_delay\_robustness\_panel.png} in both rounds.)

\emph{Trend match:} \textbf{NO} for ranking.
This mismatch is expected and \emph{does not undermine} the main
findings.  Delay remains a secondary metric in the evaluation
framework, and its Round~3 behaviour is consistent with the
conditioning-bias explanation from Section~4.5.

% ---- Best and Worst Protocol Check ------------------------------------------
\paragraph{Best and weakest protocol stability.}

\emph{Does the Round~2 best protocol remain best in Round~3?}
\textbf{No.}  \texttt{ugv\_p\_edf}, the Round~2 leader in PDR (0.630),
charge success (48.6\%), and lowest depletions (20.0), drops to rank~2
in PDR (0.732) and rank~3 in success rate (0.576) in Round~3.  It is
overtaken by \texttt{ugv\_p\_role\_priority}, which benefits most from
the removal of weather-driven routing drops.

\emph{Does the Round~2 weak tail remain weak?}
\textbf{Partially.}  \texttt{ugv\_role\_priority} remains the weakest
in charge success rate (0.449) and has the highest depletion count
(25.0) in Round~3, preserving its position at the bottom of the
survivability ranking.  However, its PDR improves enough (0.722) to
move out of last place; \texttt{ugv\_dynamic} (0.667) takes over as
the lowest-PDR protocol.

\emph{Plausible explanation for the reversal of
\texttt{ugv\_p\_role\_priority}.}
In Round~2, the CH-priority paradox (Section~4.4.4) described how
role-based protocols suffered from simultaneous CH withdrawal, which
compounded weather-driven member starvation.  In Round~3, the removal
of weather drops means that members retain connectivity to the UGV
scheduler even when both CHs charge concurrently.  This neutralises
the worst aspect of the CH-synchronisation failure mode, allowing
the preemptive role-priority policy to translate its aggressive CH
servicing into an advantage rather than a liability.  The finding
therefore does not invalidate the CH-priority paradox; instead, it
shows that the paradox is \emph{amplified by weather stress} and can
be mitigated when network conditions are favourable.

% ---- Energy per Session: the one partial match ------------------------------
\paragraph{Energy per session: a stable metric.}

Energy recovered per charging session is the only metric with a
\textbf{PARTIAL} match (Kendall $\tau = 0.81$).
\texttt{ugv\_role\_priority} leads in both rounds ($\sim$65--64\,Wh),
and the overall ranking is largely preserved with only minor
mid-tier reordering.  This stability is expected: per-session energy
depends primarily on dock hardware and session duration, neither of
which changes between rounds.
(Figures: \texttt{02\_energy\_recovered\_merged.png},
\texttt{04\_policy\_radar\_merged.png} in both rounds.)

% ---- Evidence Map -----------------------------------------------------------
\begin{table}[htbp]
\centering
\caption{Evidence map: claims, supporting figures, metric definitions,
and caveats for Section~4.6.}
\label{tab:evidence-map}
\footnotesize
\begin{tabular}{p{3.2cm} p{2.8cm} p{2.8cm} p{2.5cm} p{2.7cm}}
\hline
\textbf{Claim} & \textbf{Round~2 figure(s)} & \textbf{Round~3 figure(s)} & \textbf{Metric} & \textbf{Caveat / Note} \\
\hline

PDR ranking reversal
& \texttt{01\_mean\_pdr\_merged.png}
& \texttt{01\_mean\_pdr\_merged.png}
& Mean PDR (delivered/generated, full run)
& Kendall $\tau = -0.24$; absolute PDR higher in R3 \\

PDR temporal stability improves
& \texttt{01\_network\_pdr\_over\_time\_merged.png}
& \texttt{01\_network\_pdr\_over\_time\_merged.png}
& Windowed PDR time-series
& Reduced variance expected under fixed weather \\

Routing drops collapse in R3
& \texttt{CL\_G\_routing\_drop\_timeseries.png}
& \texttt{CL\_G\_routing\_drop\_timeseries.png}
& Charge-request ROUTING\_DROP rate
& Confirms WEATHER\_DROP was dominant cause \\

Charge success ranking shifts
& \texttt{02\_charge\_success\_rate\_merged.png}; \texttt{03\_charge\_outcome\_breakdown\_merged.png}
& \texttt{02\_charge\_success\_rate\_merged.png}; \texttt{03\_charge\_outcome\_breakdown\_merged.png}
& Charge success rate (\#STARTED / \#requests)
& $\tau = 0.33$; timeouts rise as queue contention replaces routing loss \\

Depletion gap compresses
& \texttt{03\_dead\_uav\_cumulative\_merged.png}; \texttt{CL\_F\_pdr\_vs\_depletions\_timeseries.png}
& \texttt{03\_dead\_uav\_cumulative\_merged.png}; \texttt{CL\_F\_pdr\_vs\_depletions\_timeseries.png}
& Total battery depletions per run
& $\tau = -0.05$; weather removal reduces cascade failures \\

Success--PDR association (qualitative)
& \texttt{CL\_B\_charging\_vs\_pdr\_scatter.png}; \texttt{CL\_E\_correlation\_bar\_chart.png}
& \texttt{CL\_B\_charging\_vs\_pdr\_scatter.png}; \texttt{CL\_E\_correlation\_bar\_chart.png}
& Pearson $r$ / Spearman $\rho$ (success vs.\ PDR)
& R3 per-run correlation not yet recomputed (see TODO) \\

Depletion--PDR association persists
& \texttt{CL\_F\_pdr\_vs\_depletions\_timeseries.png}; \texttt{CL\_D\_mechanism\_chain.png}
& \texttt{CL\_F\_pdr\_vs\_depletions\_timeseries.png}; \texttt{CL\_D\_mechanism\_chain.png}
& Depletions vs.\ PDR
& Compressed range weakens stat.\ power \\

Delay ranking reversal
& \texttt{06\_mean\_e2e\_delay\_merged.png}; \texttt{CL\_C\_e2e\_delay\_conditioning\_analysis.png}
& \texttt{06\_mean\_e2e\_delay\_merged.png}; \texttt{CL\_C\_e2e\_delay\_conditioning\_analysis.png}
& Mean E2E delay (delivered packets only)
& Conditioning bias; $\tau = -0.33$; delay is secondary metric \\

Energy/session ranking stable
& \texttt{02\_energy\_recovered\_merged.png}; \texttt{04\_policy\_radar\_merged.png}
& \texttt{02\_energy\_recovered\_merged.png}; \texttt{04\_policy\_radar\_merged.png}
& Mean Wh per charge session
& $\tau = 0.81$ (PARTIAL match); hardware-dominated \\

CH-priority paradox attenuated
& \texttt{CL\_M\_role\_scheduling\_audit.png}; \texttt{CL\_N\_ch\_sync\_cascade.png}
& \texttt{CL\_M\_role\_scheduling\_audit.png}; \texttt{CL\_N\_ch\_sync\_cascade.png}
& Role-stratified success rate, CH gap
& R3 members retain UGV access even during dual-CH charging \\

\hline
\end{tabular}
\end{table}

% -----------------------------------------------------------------------------
\subsubsection{Interpretation}
\label{sec:results:controlled-validation:interpretation}

\paragraph{What the validation strengthens.}
The controlled round confirms two structural properties of the
cross-layer mechanism.
First, the near-elimination of ROUTING\_DROP in Round~3 (rates of
0.010--0.020 versus 0.218--0.346 in Round~2) validates that
weather-driven routing failure was the dominant cause of request loss
in Round~2.  This confirms the feedback loop proposed in
Section~4.4.2: degraded weather $\to$ WEATHER\_DROP $\to$
ROUTING\_DROP of charge requests $\to$ missed charging $\to$ UAV
depletion $\to$ further network degradation.  When the weather trigger
is removed, the loop is broken at its entry point, and all protocols
improve.
Second, the qualitative direction of the charging--survivability--PDR
association persists: protocols that achieve higher charge success
rates tend to show fewer depletions and higher PDR, even though the
specific protocol rankings differ.  The cross-layer mechanism is
therefore not an artefact of weather variance.

\paragraph{What the validation does not prove.}
This section does not establish strict causal claims.  Both Round~2
and Round~3 are observational simulation studies, not controlled
experiments with randomised protocol assignment at the run level.
Although the design holds all factors constant except the scheduling
policy (and, across rounds, the weather mode), confounding from
residual stochastic factors cannot be fully excluded.
The fleet size (five UAVs, one UGV) is also small, limiting
generalisability to larger deployments.  Finally, the simulator
abstracts several physical-layer effects (e.g.\ exact RF propagation,
wind-induced UAV drift), so the magnitudes reported here should not
be taken as field-deployable predictions.

\paragraph{Explaining the mismatches.}
The widespread ranking reversals (9 out of 10 metrics show NO full
match) are not unexpected given the severity of the experimental
intervention.  By removing the WEATHER\_DROP channel, Round~3
eliminates the mechanism that most differentiated protocol performance
in Round~2.  When all protocols achieve routing-drop rates below 2\%,
the scheduling policy's ability to handle requests \emph{that arrive}
matters less than its ability to manage dock contention and queue
pressure---a regime in which timeout rate and waiting time become the
binding constraints.  This shift explains why protocols that were
penalised in Round~2 for poor request reachability (e.g.\
\texttt{ugv\_edf}, \texttt{ugv\_p\_role\_priority}) recover
disproportionately in Round~3.

The sole partial match---energy per session ($\tau = 0.81$)---is
consistent with the expectation that per-session energy is governed by
dock hardware parameters rather than by the scheduling or weather
layer, and therefore remains invariant across rounds.

\paragraph{Transition to Chapter~5.}
Taken together, the Round~2 and Round~3 results provide
complementary views of the same system.  Round~2 reveals the full
cross-layer mechanism under realistic weather stress, identifying
which protocols are most resilient when the feedback loop is active.
Round~3 isolates the scheduling layer by removing the environmental
trigger, showing that weather variance was a necessary amplifier of
the performance gaps but not the sole source of protocol
differentiation.  Both rounds agree that survivability---measured by
battery depletions---remains the central mediator between charging
outcomes and network quality.  Chapter~5 synthesises these findings,
discusses their implications for scheduling-policy design in
shared-substrate FANET systems, and outlines directions for future
work.

\label{sec:results:test-round3}
